\section{Task Scheduling Problem}
\label{sec:g-task-scheduling}
\problemlink{https://atcoder.jp/contests/abc103/tasks/abc103_a}{AtCoder ABC103 A}
\source{AtCoder ABC103 A (Notion: Greedy)}{Easy}

\problemstatement{There are three tasks with values $A_1,A_2,A_3$
($1\le A_i\le 100$). The first task you complete costs $0$; completing task
$j$ right after task $i$ costs $|A_j-A_i|$. Find the minimum total cost to
complete all three tasks.}

\begin{patternbox}
``Visit points on a line, pay the distance between consecutive visits'' ---
the cost of any route is at least the length of the segment it covers.
\end{patternbox}

\subsection*{Approach and intuition}

Put the three values on a number line. Any order of visiting them must
travel from the smallest to the largest value at some point, so every order
costs at least $\max A-\min A$. Visiting them in sorted order walks
monotonically from $\min$ to $\max$ and costs exactly that amount.

\begin{ideabox}[Why It Is Optimal]
\emph{Lower bound:} the route touches both $\min A$ and $\max A$; by the
triangle inequality the total path length between those two visits is at
least $\max A-\min A$. \emph{Achievability:} sorted order has cost
$(A_{(2)}-A_{(1)})+(A_{(3)}-A_{(2)})=A_{(3)}-A_{(1)}$.
\end{ideabox}

\subsection*{Algorithm}
\begin{algorithmic}[1]
\State read $A_1,A_2,A_3$
\State \Return $\max(A_1,A_2,A_3)-\min(A_1,A_2,A_3)$
\end{algorithmic}

\dryrun{$A=(1,6,3)$: sorted $1,3,6$, cost $2+3=5=6-1$. \checkmark}

\subsection*{C++17 implementation}
\cppfile{greedy/01_task_scheduling.cpp}

\begin{complexitybox}
\textbf{Time:} $O(1)$. \textbf{Space:} $O(1)$.
\end{complexitybox}

\begin{pitfallbox}
\begin{itemize}
  \item All equal values ($100\ 100\ 100$) give $0$ --- the formula
        handles it.
  \item The same idea generalises to $n$ points: answer $\max-\min$, not a
        sum of sorted gaps you must compute one by one.
\end{itemize}
\end{pitfallbox}
